\section{结论}

本文首先构建了一个多模态数据集 ChartRWC，其中每条数据均由图表及其对应文本描述组成。文本部分主要用于解释或解读图表内容，为自然语言输入驱动的图表类型推荐任务提供支撑材料。

在方法设计方面，本文采用两种互补的检索策略。第一，使用基于关键词匹配的倒排索引方法，快速定位与输入文本直接相关的候选图表类型。第二，使用基于向量检索的语义匹配方法，将输入文本和数据集中图表的文本描述分别编码为向量表示，再通过计算向量相似度获取与输入最接近的候选图表类型。

在推荐阶段，两种方法分别生成候选结果，随后通过基于排序的融合策略对两个结果列表进行加权整合，从而兼顾推荐准确性与多样性。在候选结果评估环节，本文引入期望倒数排名指标（Expected Reciprocal Rank, ERR）对融合结果进行性能评价。

在图表生成环节，系统利用 ECharts 库将推荐结果转化为可视化输出，实现从文本输入到图表的转换。案例研究和用户研究证明了本文方法的有用性与有效性。

\newpage
